\section{Dimension induction}
\label{sec:induction}

We first prove the desired inequality in the range reached by dimension
induction.  Every decimal in this section denotes an exact rational number.

\begin{theorem}\label{thm:induction}
Let $f:\{-1,1\}^n\to\{-1,1\}$, let $m=\E f$, and let
$0<\rho<1$.  If $s=\sqrt{1-\rho^2}\le .84$, then
\[
  \E H(T_\rho f)\ge H(m)-1+s.
\]
\end{theorem}

We strengthen the desired bound to an inequality for a nonsymmetric
function of the noisy output.  Averaging this inequality with its
reflected version will recover $H$.
The induction uses an asymmetric function $G$ and a polynomial profile $B$.
The scalar inequalities that support it are stated in Lemma~\ref{lem:scalar}
and proved by the finite arithmetic certificates in Appendix~\ref{app:scalar}.
We first establish the two reductions that explain their domains.

\subsection{Compression and a small coordinate}

The following sorting argument is the concave-profile form of the
compression principle of Anantharam, Bogdanov, Chakrabarti, Jayram, and
Nair~\cite[Proposition~2]{ABCJN}. We include the short proof to specify
its direction for $H$.
We order the cube coordinatewise using $-1<1$. A function is
\emph{increasing} if it is nondecreasing in each coordinate.

\begin{lemma}[Compression]\label{ind:lem:compression}
For every Boolean function $f$ there is an increasing Boolean function
$\widetilde f$ with the same mean such that
\[
  \E H(T_\rho\widetilde f)\le \E H(T_\rho f).
\]
\end{lemma}

\begin{proof}
Fix a coordinate, and write $f_-,f_+$ for its two sections.  Replace them
pointwise by $\min(f_-,f_+),\max(f_-,f_+)$, respectively.  Let $T$ denote
noise in the remaining coordinates and put
\[
 a=T\bigl((f_-+f_+)/2\bigr),\qquad
 c=T\bigl((f_+-f_-)/2\bigr).
\]
Sorting preserves $a$ and replaces $c$ by
\[
 c'=T\bigl(|f_+-f_-|/2\bigr)\ge |c|,
\]
since $T$ has a nonnegative kernel.  The two fully noised values are
$a\pm\rho c$ before sorting and $a\pm\rho c'$ afterwards.  For fixed $a$,
the function $u\mapsto H(a+u)+H(a-u)$ is even and concave on its domain,
and hence nonincreasing for $u\ge0$.  Thus sorting weakly decreases the
pair average of $H$.  It preserves the mean.

Perform this operation once in each coordinate.  If the function is
already increasing in a different coordinate, both its sections are
increasing in that coordinate; their pointwise minimum and maximum remain
increasing. The resulting function is therefore increasing in every
coordinate.
\end{proof}

Call a coordinate \emph{active} if the function depends on it.  Inactive
coordinates may be omitted from all expectations below.  Define
\begin{equation}\label{ind:eq:b}
 b(m)=\min\left\{\frac38,\frac9{16}-\frac{|m|}{2},1-|m|\right\},
 \qquad -1\le m\le1.
\end{equation}

\begin{lemma}[Choice of a coordinate]\label{ind:lem:coordinate}
If an increasing Boolean function of mean $m$ has at least four active
coordinates, one of them has first-level coefficient
\[
 0<x=\E[f(X)X_i]\le b(m).
\]
The two sections in this coordinate have means $m-x$ and $m+x$.
\end{lemma}

\begin{proof}
For an increasing Boolean function, the coefficient of coordinate $i$ is
\[
 x_i=\E\bigl[(f_+-f_-)/2\bigr].
\]
The integrand takes values in $\{0,1\}$, so $x_i>0$ precisely when the
coordinate is active.  Choose a coordinate with the least active
coefficient $x$.  For any four distinct active coordinates,
\[
 x\le\frac{x_{i_1}+x_{i_2}+x_{i_3}+x_{i_4}}4
   =\E[f(X)Z],\qquad
 Z=\frac{X_{i_1}+X_{i_2}+X_{i_3}+X_{i_4}}4.
\]
The distribution of $Z$, in decreasing order, is
\[
\begin{array}{c|rrrrr}
z&1&1/2&0&-1/2&-1\\ \hline
\Pr(Z=z)&1/16&4/16&6/16&4/16&1/16.
\end{array}
\]
At fixed mean $m$, the largest possible $\E[fZ]$ is obtained by placing
the positive values of $f$ at the largest values of $Z$, allowing a
fraction of the threshold level.  Indeed, transfer of positive mass from
a smaller value of $Z$ to a larger one can only increase the objective;
allowing fractional values is a relaxation and therefore gives an upper
bound for Boolean $f$.  Since $\E Z=0$ and the law of $Z$ is symmetric,
this maximum is twice the integral over its top fraction
$q=(1-|m|)/2$.  For $q\in[0,1/2]$ that quantity is
\[
 \begin{cases}
 2q,&0\le q\le1/16,\\
 q+1/16,&1/16\le q\le5/16,\\
 3/8,&5/16\le q\le1/2.
 \end{cases}
\]
These are exactly the three pieces of $b(m)$.  Finally, the section
means have average $m$ and half-difference $x$, proving the last claim.
\end{proof}

\subsection{The scalar interface}

For the fixed value of $s$, use one row of
Table~\ref{ind:tab:parameters} whose interval contains $s$.  The profile
polynomials use a finer subdivision of these same intervals, specified
in Appendix~\ref{app:scalar}.  At a shared endpoint, choose the parameters
and the profile consistently from one adjoining interval.  Put
$\delta=(1-\rho)/2$, so that $4\delta(1-\delta)=s^2$; the auxiliary
bounds in the table satisfy $r_0^2\le\delta\le r_1^2$.

\begin{table}[htbp]
\centering
\small
\setlength{\tabcolsep}{3.5pt}
\begin{tabular}{rrrrrrrrrrr}
\toprule
$i$ & $s_{\rm lo}$ & $s_{\rm hi}$ & $r_0$ & $r_1$ & $L$ & $U$ & $V$
    & $\alpha$ & $\Lambda$ & $\mu$\\
\midrule
0 & 0   & 70  & 0   & 36  & 940 & $-440$ & 300 & 960 & 1000 & 956\\
1 & 70  & 160 & 35  & 81  & 940 & $-440$ & 300 & 930 & 1050 & 953\\
2 & 160 & 300 & 80  & 152 & 886 & $-398$ & 250 & 860 & 1150 & 985\\
3 & 300 & 490 & 151 & 254 & 842 & $-335$ & 180 & 770 & 1150 & 1024\\
4 & 490 & 630 & 253 & 335 & 842 & $-335$ & 180 & 690 & 1150 & 1058\\
5 & 630 & 750 & 334 & 412 & 842 & $-335$ & 180 & 632 & 1160 & 1087\\
6 & 750 & 840 & 411 & 479 & 842 & $-335$ & 180 & 583 & 1160 & 1111\\
\bottomrule
\end{tabular}
\caption{Induction parameters.  Every displayed entry except $i$ is
$1000$ times the corresponding parameter.}
\label{ind:tab:parameters}
\end{table}

Define
\begin{equation}\label{ind:eq:G}
 \begin{split}
 \lambda(h)&=L+Uh+Vh^2,\\
 G(p)&=H(2p-1)\left(1+
 \frac{2p-1}{1+H(2p-1)\lambda(H(2p-1))}\right),\qquad 0\le p\le1.
 \end{split}
\end{equation}
All rows have $U<0$, $V\ge0$, and $L+U\ge .488$.  Thus
$\lambda(h)\ge L+U>0$ for $0\le h\le1$, and $G$ is nonnegative,
with $G(0)=G(1)=0$.  Its defining asymmetry cancels under reflection:
\begin{equation}\label{ind:eq:G-reflection}
 G(p)+G(1-p)=2H(2p-1).
\end{equation}

The profiles in the next lemma are explicit polynomials of the form
\begin{equation}\label{ind:eq:profile-def}
 B(m)=(1-m^2)P(m),\qquad
 P(m)=1+10^{-8}\sum_{j\ge1}c_j\bigl(T_j(m)-T_j(0)\bigr),
\end{equation}
where only finitely many $c_j$ are nonzero, and
$T_0(u)=1$, $T_1(u)=u$, $T_{j+1}(u)=2uT_j(u)-T_{j-1}(u)$.
The integer coefficients and their arithmetic verification are given in
Appendix~\ref{app:scalar}.

With these functions, the estimate we will propagate is
\[
 \E G\bigl((1+T_\rho f)/2\bigr)\ge sB(m)
 \qquad(m=\E f)
\]
for every increasing Boolean function $f$.
Its dual $f^*(x)=-f(-x)$ is increasing and has mean $-m$.
Noise commutes with input reversal, so averaging the proposed estimates
for $f$ and $f^*$ and using Equation~\eqref{ind:eq:G-reflection} would give
\[
 \E H(T_\rho f)\ge\frac{s}{2}\bigl(B(m)+B(-m)\bigr).
\]
Thus we need a profile whose reflected average reaches $H(m)-1+s$,
and an estimate for $G$ that passes from two sections to their parent.
The next lemma supplies the two-point gap used in this passage, the
reflected comparison, the condition that closes induction on the
small-coordinate domain, and the initial cases.

\begin{lemma}[Scalar inequalities]\label{lem:scalar}
For every $0<s\le .84$, choose the parameters and the polynomial $B$
as above.  Then $B\ge0$ on $[-1,1]$, $B(0)=1$, $B(\pm1)=0$, and
the following four assertions hold.
\begin{enumerate}
\item For $0\le p_0\le p_1\le1$, write $z=p_1-p_0$ and
$p_u=p_0+uz$, and set
\[
 D=\frac{G(p_\delta)+G(p_{1-\delta})-G(p_0)-G(p_1)}2,
 \qquad S=\alpha G(p_1)+(1-\alpha)G(p_0).
\]
Then $D\ge0$ and
\begin{equation}\label{ind:eq:pair}
 D(S+\Lambda D)\ge
 \mu s^2\left(z-\frac{G(p_1)-G(p_0)}{4L}\right)^2.
\end{equation}
\item For every $m\in[-1,1]$,
\begin{equation}\label{ind:eq:reflection-profile}
 \frac{s}{2}\bigl(B(m)+B(-m)\bigr)\ge H(m)-1+s.
\end{equation}
\item For $m\in[-1,1]$ and $0<x\le b(m)$, put
\[
 B_\pm=B(m\pm x),\qquad
 C=\frac{B_+-B_-}{2x},\qquad
 A=\frac{2B(m)-B_+-B_-}{2x^2}.
\]
Then
\begin{equation}\label{ind:eq:profile}
 A\bigl[B(m)+(2\alpha-1)xC+(\Lambda-1)x^2A\bigr]
 \le \mu\left(1-\frac{s}{2L}C\right)^2.
\end{equation}
\item Every increasing Boolean function depending on at most three
coordinates satisfies
\begin{equation}\label{ind:eq:base}
 \E G\bigl((1+T_\rho f)/2\bigr)\ge sB(\E f).
\end{equation}
\end{enumerate}
\end{lemma}

The four assertions of Lemma~\ref{lem:scalar}, including the entire
three-coordinate base case, are proved in Appendix~\ref{app:scalar}.
We now show how the pair bound and profile comparison propagate the
asymmetric estimate from three active coordinates to arbitrary dimension.

\subsection{Averaging and dimension induction}

\begin{proposition}\label{ind:prop:asymmetric}
For the fixed $s$, parameters, and profile above, every increasing
Boolean function satisfies
\begin{equation}\label{ind:eq:all-dimensions}
 \E G\bigl((1+T_\rho f)/2\bigr)\ge sB(\E f).
\end{equation}
\end{proposition}

\begin{proof}
Induct on the number of active coordinates.  Lemma~\ref{lem:scalar}(4)
provides the cases with at most three.  Otherwise choose the coordinate
in Lemma~\ref{ind:lem:coordinate}, and denote its increasing sections by
$f_-$ and $f_+$, with means $m-x$ and $m+x$.  On the smaller cube let
\[
 p_0=\frac{1+T_\rho f_-}{2},\qquad
 p_1=\frac{1+T_\rho f_+}{2}.
\]
In these expressions, $T_\rho$ acts only on the remaining coordinates.  Positivity of its kernel gives $p_0\le p_1$, and
$\E(p_1-p_0)=x$.  At the two values of the chosen coordinate the parent
noise probabilities are $p_\delta$ and $p_{1-\delta}$.

Apply Equation~\eqref{ind:eq:pair} pointwise, allowing $p_0=p_1$, in
which case both $D$ and the expression on the right vanish.  Set
\[
 y=\frac{\E G(p_1)}s,\qquad t=\frac{\E G(p_0)}s,\qquad
 e=\frac{\E D}s,\qquad k=\frac{s}{2L},\qquad
 Q=\alpha y+(1-\alpha)t.
\]
All of $y,t,e,Q$ are nonnegative.  To average the pair inequality, write
$w=p_1-p_0-(G(p_1)-G(p_0))/(4L)$.  Cauchy--Schwarz yields
\[
 \begin{split}
 (\E D)\,\E(S+\Lambda D)
 &\ge\left(\E\sqrt{D(S+\Lambda D)}\right)^2\\
 &\ge\mu s^2(\E|w|)^2
 \ge\mu s^2(\E w)^2.
 \end{split}
\]
No fixed sign of $w$ is required.  Since
$\E w=x-k(y-t)/2$, division by $s^2$ gives
\begin{equation}\label{ind:eq:averaged-pair}
 e(Q+\Lambda e)\ge\mu\left[x-\frac{k}{2}(y-t)\right]^2.
\end{equation}
The parent expectation divided by $s$ equals $(y+t)/2+e$.  Solving
Equation~\eqref{ind:eq:averaged-pair} for the nonnegative variable $e$
therefore gives the lower bound
\begin{equation}\label{ind:eq:R}
 R_x(y,t)=\frac{y+t}{2}+
 \frac{\sqrt{Q^2+4\Lambda\mu v^2}-Q}{2\Lambda},
 \qquad v=x-\frac{k}{2}(y-t).
\end{equation}

We next justify substitution of the induction bounds into this root.
Put $\Delta=\sqrt{Q^2+4\Lambda\mu v^2}$.  For $y,t>0$,
\[
 \begin{split}
 \partial_yR_x
 &=\frac12-\frac{\alpha}{2\Lambda}
    +\frac{\alpha Q-2\Lambda\mu kv}{2\Lambda\Delta},\\
 \partial_tR_x
 &=\frac12-\frac{1-\alpha}{2\Lambda}
    +\frac{(1-\alpha)Q+2\Lambda\mu kv}{2\Lambda\Delta}.
 \end{split}
\]
Using $Q\ge0$ and
$|v|/\Delta\le1/(2\sqrt{\Lambda\mu})$, each derivative is at least
\begin{equation}\label{ind:eq:root-derivative}
 \frac12\left(1-\frac{\max(\alpha,1-\alpha)}{\Lambda}
              -k\sqrt{\frac\mu\Lambda}\right)>0.
\end{equation}
For completeness, the strict sign follows by rational arithmetic from
Table~\ref{ind:tab:parameters}.  If
$a=1-\max(\alpha,1-\alpha)/\Lambda$, every row has $a>0$ and
\[
 a^2-\left(\frac{s_{\rm hi}}{2L}\right)^2\frac\mu\Lambda
 \ge\frac{12133}{44180000}>0.
\]
Since $s\le s_{\rm hi}$, this proves the sign in
Equation~\eqref{ind:eq:root-derivative}.  Continuity extends
coordinatewise monotonicity of $R_x$ to the whole nonnegative quadrant.

Each section has fewer active coordinates, so the induction hypothesis
implies $y\ge B_+$ and $t\ge B_-$.  We have consequently obtained the
lower bound $R_x(B_+,B_-)$ for the normalized parent expectation.
In the notation of Lemma~\ref{lem:scalar}(3), at these arguments
\[
 \frac{B_++B_-}{2}=B(m)-x^2A,\qquad
 Q=B(m)-x^2A+(2\alpha-1)xC,\qquad
 v=x(1-kC).
\]
If $A\le0$, the first of these terms is already at least $B(m)$,
and the root correction in Equation~\eqref{ind:eq:R} is nonnegative.
Suppose $A>0$.  The correction is the nonnegative solution $u$ of
$u(Q+\Lambda u)=\mu x^2(1-kC)^2$.  Because $Q\ge0$ and $\Lambda>0$,
the left side is increasing for $u\ge0$.  Equation~\eqref{ind:eq:profile},
multiplied by $x^2$, states that
\[
 \begin{split}
 (x^2A)(Q+\Lambda x^2A)
 &=x^2A\bigl[B(m)+(2\alpha-1)xC+(\Lambda-1)x^2A\bigr]\\
 &\le\mu x^2(1-kC)^2.
 \end{split}
\]
Thus the correction is at least $x^2A$, and again
$R_x(B_+,B_-)\ge B(m)$.  This completes the induction.
\end{proof}

\begin{proof}[Proof of Theorem~\ref{thm:induction}]
By Lemma~\ref{ind:lem:compression}, it suffices to consider increasing
$f$.  Its dual $f^*(x)=-f(-x)$ is also increasing and has mean $-m$.
Noise commutes with input reversal, so
$T_\rho f^*(x)=-T_\rho f(-x)$.  Apply
Proposition~\ref{ind:prop:asymmetric} to $f$ and $f^*$, average the two
inequalities, and change variables $x\mapsto-x$ in the second
expectation.  Equation~\eqref{ind:eq:G-reflection} gives
\[
 \E H(T_\rho f)
 \ge\frac{s}{2}\bigl(B(m)+B(-m)\bigr)
 \ge H(m)-1+s,
\]
where the last inequality is Equation~\eqref{ind:eq:reflection-profile}.
\end{proof}
